% ===========================================================================
\section{Boundary functions of discontinuous functions}\label{sec:discont}
% ===========================================================================
Sources: \cite[Chapter II]{Diss}; \cite[Theorems 4--6]{JMM}. We work in $\Hp$ with boundary $\R$.

\subsection{Baire functions}

\begin{definition}\label{def:meets-in}
\lean{TJK.Discont.MeetsIn}
\uses{def:arcs}
For arcs $\gamma_x,\gamma_y$ at $x,y\in\R$ and $c>0$, $\gamma_x$ \emph{meets $\gamma_y$ below height $c$} if there are
subarcs $\gamma_x'\subseteq\gamma_x$ at $x$ and $\gamma_y'\subseteq\gamma_y$ at $y$, contained in $\{\operatorname{Im}z\le c\}$,
with $\gamma_x'\cap\gamma_y'\ne\emptyset$. This relation is symmetric, and if it holds for $c$ it holds for every $c'\ge c$.
\end{definition}

\begin{lemma}[{the ``Figure 5'' argument of \cite[proof of Theorem 8]{Diss}}]\label{lem:figure5}
\lean{TJK.Discont.figure5}
\uses{lem:crosscut-H, fact:tietze-arc}
Let $x<t<y$ be real, let $\alpha,\beta$ be arcs at $x,y$ in $\Hp$, contained in $\{\operatorname{Im}z\le c\}$, with
$\alpha\cap\beta\ne\emptyset$, and let $\tau$ be an arc at $t$ in $\Hp$ whose other endpoint $p$ has $\operatorname{Im}p=c$.
Then $\tau$ meets $\alpha\cup\beta$.
\end{lemma}
\begin{proof}
$\alpha\cup\beta$ is path connected, so by Fact~\ref{fact:tietze-arc} it contains a simple arc $\sigma$ from $x$ to $y$, a
crosscut of $\Hp$ with $\max_\sigma\operatorname{Im}\le c$. Suppose $\tau\cap\sigma=\emptyset$, and let $V,U$ be the bounded and
unbounded components of $\Hp\setminus\sigma$ (Lemma~\ref{lem:crosscut-H}). The connected set $\tau\setminus\{t\}$ lies in
$U\cup V$, hence in $U$ or in $V$; and $t\in\overline{\tau\setminus\{t\}}$. If it lies in $U$, then
$t\in\overline U\cap\R=\R\setminus(x,y)$, which is false. If it lies in $V$, then $p\in V$, so $p+i\eps\in V$ for small $\eps>0$;
but $p+i\eps\in U$ since it lies above $\sigma$. Contradiction.
\end{proof}

\begin{theorem}[{\cite[Theorem 8]{Diss}}; {\cite[Theorem 4]{JMM}}]\label{thm:baire-plus-one}
\lean{TJK.Discont.baireClass_succ_of_boundaryFn}
\uses{def:meets-in, lem:figure5, lem:strip, lem:subarc, lem:intervals-open, lem:equiv, lem:monotone-extend,
  lem:monotone-compose, lem:baire-precomp, lem:baire-imp-class, lem:additive-basic, lem:baire-criterion}
Let $Y$ be a separable arcwise connected metric space, $\xi\ge1$, $f:\Hp\to Y$ of Baire class $\xi$, $E\subseteq\R$, and
$\varphi:E\to Y$ a boundary function for $f$. Then $\varphi$ is of Baire class $\xi+1$.
\end{theorem}
\begin{proof}
Let $U\subseteq Y$ be open, $V=Y\setminus\overline U$, $A=\varphi^{-1}(U)$, $B=\varphi^{-1}(V)$, $P=A\cup B$. By
Lemma~\ref{lem:baire-criterion} (with $\xi+1\ge2$) it suffices to find $T\in P^{\xi+2}(\R)$ with $A\subseteq T$ and
$T\cap B=\emptyset$: then $\varphi^{-1}(U)\subseteq T\cap E\subseteq\varphi^{-1}(\overline U)$ and $T\cap E\in P^{\xi+2}(E)$.
For $x\in P$ choose (Lemma~\ref{lem:subarc}) an arc $\gamma_x$ at $x$ with $f\to\varphi(x)$ along it,
$\gamma_x\subseteq\{\abs{z-x}\le1\}$, and $f(\gamma_x\setminus\{x\})\subseteq U$ if $x\in A$, $\subseteq V$ if $x\in B$. Then
$\gamma_x\cap\gamma_y=\emptyset$ for $x\in A$, $y\in B$. Write $c_n=1/(n+1)$ and ``meets at $n$'' for ``meets below height
$c_n$''. Let
\[
L=\{x\in P:\forall n\ \exists y\in P\setminus\{x\},\ \gamma_x\text{ meets }\gamma_y\text{ at }n\},\qquad M=P\setminus L,
\]
and let $K_n$ be the set of $x\in P$ such that $\gamma_x$ reaches the line $\operatorname{Im}z=c_n$ and, if $x\in M$, $\gamma_x$
meets no $\gamma_y$, $y\in P\setminus\{x\}$, at $n$. Every $x\in P$ is in some $K_n$, and every $x\in M$ is in $K_n$ for all
large $n$ (an arc at $x$ reaches every small height; and ``meets'' is monotone in the height).

\emph{Step 1: $L\cap A\approx G_a\cap P$ and $L\cap B\approx G_b\cap P$ for $\Gd$ sets $G_a,G_b$.} Let $x\in L\cap A$,
$y\in P\setminus\{x\}$ with $\gamma_x$ meeting $\gamma_y$ at $n$. We claim $[x,y]\cap K_n\subseteq L\cap A$ (with $[x,y]$ the
closed interval between $x$ and $y$). First, $\gamma_x\cap\gamma_y\ne\emptyset$, so $y\in A$. Let $t\in[x,y]\cap K_n$, $t\ne x$.
If $t=y$, $\gamma_t$ meets $\gamma_x$ at $n$. Otherwise let $\tau$ be the subarc of $\gamma_t$ from $t$ to its first point on
$\operatorname{Im}z=c_n$; by Lemma~\ref{lem:figure5}, $\tau$ meets one of the two meeting subarcs, so $\gamma_t$ meets $\gamma_x$
or $\gamma_y$ at $n$. Since $t\in K_n$, this gives $t\notin M$, so $t\in L$; and $\gamma_t$ meets an arc of a point of $A$, so
$t\in A$. Let $W_n$ be the union of the intervals $[x,y]$ over all such pairs $(x,y)$. Then
$L\cap A=\bigcap_n(W_n\cap P)$: ``$\subseteq$'' since $x\in[x,y]$; and if $t\in P\cap\bigcap_nW_n$, take $m$ with $t\in K_m$, then
$t\in L\cap A$ by the claim. By Lemma~\ref{lem:intervals-open} each $W_n\approx O_n$ with $O_n$ open; let $G_a=\bigcap_nO_n$. By
Lemma~\ref{lem:equiv}, $L\cap A\approx G_a\cap P$. Exchanging $A$ and $B$ gives $G_b$.

\emph{Step 2: $M\cap A\approx D\cap M$ with $D\in P^{\xi+2}(\R)$.} For $x\in K_n$ let $p_n(x)$ be the first point of
$\operatorname{Im}z=c_n$ reached along $\gamma_x$ from $x$. If $x<x'$ in $M\cap K_n$, the subarcs from $x$ to $p_n(x)$ and from
$x'$ to $p_n(x')$ are disjoint (they would otherwise meet at $n$), so $\operatorname{Re}p_n(x)<\operatorname{Re}p_n(x')$ by
Lemma~\ref{lem:strip}. Also $\abs{x-\operatorname{Re}p_n(x)}\le1$. By Lemma~\ref{lem:monotone-extend} there is a weakly increasing
$h_n:\R\to\R$ equal to $\operatorname{Re}p_n$ on $M\cap K_n$; let $g_n(x)=f(h_n(x)+ic_n)$. By Lemmas~\ref{lem:baire-precomp}
and~\ref{lem:monotone-compose} there is a countable $N$ such that each $g_n|_{M\setminus N}$ is of Baire class $\xi$. For
$x\in M$ and $n$ large, $g_n(x)=f(p_n(x))$ with $p_n(x)\in\gamma_x\setminus\{x\}$ and $\operatorname{Im}p_n(x)\to0$; as $\gamma_x$
is compact and meets $\R$ only at $x$, $p_n(x)\to x$ along $\gamma_x$, so $g_n(x)\to\varphi(x)$. Hence $\varphi|_{M\setminus N}$ is
of Baire class $\xi+1$, hence of class $P^{\xi+2}$ (Lemma~\ref{lem:baire-imp-class}), so $A\cap(M\setminus N)=D\cap(M\setminus N)$
for some $D\in P^{\xi+2}(\R)$ (Lemma~\ref{lem:additive-basic}(f)).

\emph{Step 3.} From Step 1, $L\approx(G_a\cup G_b)\cap P$, so
$M\cap A\approx D\cap M\approx D\cap\bigl(P\setminus(G_a\cup G_b)\bigr)$ and
\[
A=(L\cap A)\cup(M\cap A)\approx T_0\cap P,\qquad T_0=G_a\cup\bigl(D\setminus(G_a\cup G_b)\bigr)\in P^{\xi+2}(\R),
\]
since $G_a\in Q^2\subseteq P^{\xi+2}$ and $\R\setminus(G_a\cup G_b)$ is $\Fs$. With the countable sets
$C_1=A\setminus T_0$ and $C_2=(T_0\cap P)\setminus A$, the set $T=(T_0\setminus C_2)\cup C_1$ is in $P^{\xi+2}(\R)$ (countable sets
are $\Fs$), contains $A$, and misses $B$.
\end{proof}

\begin{corollary}[{\cite[Theorem 8, Corollary]{Diss}}; {\cite[Theorems 4, 5]{JMM}}]\label{cor:borel}
\lean{TJK.Discont.measurable_of_boundaryFn, TJK.Discont.baireClass_succ_of_boundaryFn_real}
\uses{thm:baire-plus-one, fact:borel-baire, lem:baire-imp-class, lem:additive-basic}
Let $Y$ be a separable, arcwise connected metric space, $f:\Hp\to Y$ Borel measurable, and $\varphi:E\to Y$ a boundary function for $f$.
Then $\varphi$ is Borel measurable. In particular (\cite[Theorems 4, 5]{JMM}), if $f$ is a real function of Baire class
$\xi\ge1$ in the upper half-plane, every finite-valued boundary function of $f$ is of Baire class $\xi+1$, and if $f$ is
Borel measurable so is every finite-valued boundary function.
\end{corollary}
\begin{proof}
By Fact~\ref{fact:borel-baire}, $f$ is of some Baire class $\xi<\omega_1$ (if $\xi=0$, also of class $1$); by
Theorem~\ref{thm:baire-plus-one}, $\varphi$ is of Baire class $\xi+1$, hence Borel measurable
(Lemma~\ref{lem:baire-imp-class}). $\R$ is separable and arcwise connected.
\end{proof}

\subsection{Lebesgue measure}
$X_0=\R\times\{0\}$, $X_1=\R\times\{1\}$; $m$ is planar Lebesgue measure, $m^1$ linear measure on a line, and $m_e$, $m^1_e$
the outer measures. A \emph{segment} below means a closed line segment with one endpoint $\langle a_0(s),0\rangle\in X_0$ and the
other $\langle a_1(s),1\rangle\in X_1$; write $\ell_s(y)=(1-y)a_0(s)+ya_1(s)$.

\begin{definition}[trapezoids]\label{def:trapezoid}
\lean{TJK.Discont.closedTrapezoidE, TJK.Discont.closedTrapezoid, TJK.Discont.IsTrapezoid, TJK.Discont.segTrap}
For $a_0\le b_0$, $a_1\le b_1$ in $[-\infty,+\infty]$ with $a_i<+\infty$ and $b_i>-\infty$ (and $0\cdot(\pm\infty)=0$), let
$T(a_0,b_0;a_1,b_1)=\{\langle x,y\rangle:0\le y\le1,\ x\in\R,\ (1-y)a_0+ya_1\le x\le(1-y)b_0+yb_1\}$, a \emph{closed trapezoid}; $\emptyset$
is also one. $S$ is a \emph{trapezoid} if $S=\emptyset$ or $T^\circ\subseteq S\subseteq T$ for a closed trapezoid $T$, where
$T^\circ=\{\langle x,y\rangle:0\le y\le1,\ (1-y)a_0+ya_1<x<(1-y)b_0+yb_1\}$ is the interior of $T$ \emph{relative to the strip}
$\{0\le y\le1\}$. For segments $s,s'$, $T(s,s')$ is the set of $\langle x,y\rangle$ with $0\le y\le1$ and $x$ between $\ell_s(y)$
and $\ell_{s'}(y)$; for disjoint $s,s'$ with $a_i(s)\le a_i(s')$ this is $T(a_0(s),a_0(s');a_1(s),a_1(s'))$, and $T(s,s)=s$.
(In Lean a segment is given by its pair of abscissas $\langle a_0(s),a_1(s)\rangle$.)
\end{definition}
A trapezoid $S$ is Lebesgue measurable and
$m(S)=\frac12\bigl[(b_0-a_0)+(b_1-a_1)\bigr]=\frac12\bigl[m^1(S\cap X_0)+m^1(S\cap X_1)\bigr]$: integrate the lengths
$(1-y)(b_0-a_0)+y(b_1-a_1)$ of the horizontal sections, and note that $S\cap X_i$ lies between $(a_i,b_i)$ and $[a_i,b_i]$.
(The restrictions $a_i<+\infty$, $b_i>-\infty$ exclude the degenerate sides $a_i=b_i=\pm\infty$, for which the formula fails.)
Only the case of finite parameters and $S=T$ is used below (\texttt{volume\_closedTrapezoid}). Two disjoint segments $s\ne s'$
satisfy either $a_i(s)<a_i(s')$ for $i=0,1$ or $a_i(s)>a_i(s')$ for $i=0,1$ (otherwise $\ell_s-\ell_{s'}$, which is affine, changes sign
on $[0,1]$), hence either $\ell_s<\ell_{s'}$ on $[0,1]$ or $\ell_s>\ell_{s'}$ on $[0,1]$. In particular a set $\Sigma$ of pairwise
disjoint segments is totally ordered by $s\preceq s'\iff a_0(s)\le a_0(s')\text{ and }a_1(s)\le a_1(s')$.

\begin{lemma}[{\cite[Lemma 23]{Diss}}]\label{lem:th-lemma23}
\lean{TJK.Discont.segTrap_subset_union}
\uses{def:trapezoid}
Let $s_1,s_2,s_3,s_4$ be segments, any two equal or disjoint. If $T(s_1,s_2)\cap T(s_3,s_4)\ne\emptyset$ then
$T(s_1,s_3)\subseteq T(s_1,s_2)\cup T(s_3,s_4)$.
\end{lemma}
\begin{proof}
The horizontal section of $T(s_i,s_j)$ at height $y$ is the interval between $\ell_{s_i}(y)$ and $\ell_{s_j}(y)$. The order of the
four numbers $\ell_{s_i}(y)$ (with ties exactly between equal segments) does not depend on $y$. Whether two such intervals meet
depends only on this order; so if the two sections meet at one height they meet at every height. The union of two meeting
intervals is an interval, and it contains $\ell_{s_1}(y)$ and $\ell_{s_3}(y)$.
\end{proof}

\begin{lemma}[{\cite[Lemma 24]{Diss}}]\label{lem:th-lemma24}
\lean{TJK.Discont.isTrapezoid_biUnion}
\uses{def:trapezoid}
If $\Sigma$ is a set of pairwise disjoint segments, $\bigcup_{s,s'\in\Sigma}T(s,s')$ is a trapezoid.
\end{lemma}
\begin{proof}
We may assume $\Sigma\ne\emptyset$. Let $A_i=\inf_sa_i(s)$, $B_i=\sup_sa_i(s)$ (in $[-\infty,+\infty]$) and $T=T(A_0,B_0;A_1,B_1)$.
Since $(1-y)A_0+yA_1\le\ell_s(y)\le(1-y)B_0+yB_1$ for all $s$, the union lies in $T$. Let $\langle x,y\rangle\in T^\circ$. From
$(1-y)A_0+yA_1<x$ we get $s_1,s_2\in\Sigma$ with $(1-y)a_0(s_1)+ya_1(s_2)<x$ (if $y=1$ or $y=0$
only one parameter matters; if $A_0=-\infty$ and $y<1$ take any $s_2$ and $a_0(s_1)$ small; symmetrically if $A_1=-\infty$ and $y>0$;
if both are finite approximate each within $\delta=\frac12[x-(1-y)A_0-yA_1]$). The
$\preceq$-smaller $s$ of $s_1,s_2$ has $\ell_s(y)\le(1-y)a_0(s_1)+ya_1(s_2)<x$. Likewise there is $s'\in\Sigma$ with $x<\ell_{s'}(y)$,
so $\langle x,y\rangle\in T(s,s')$.
\end{proof}

\begin{theorem}[{\cite[Theorem 9]{Diss}}]\label{thm:th-thm9}
\lean{TJK.Discont.outerMeasure_union_segments, TJK.Discont.outerMeasure_union_segments_le, TJK.Discont.le_outerMeasure_union_segments}
\uses{def:trapezoid}
Let $\Sigma$ be a set of pairwise disjoint segments and $S=\bigcup\Sigma$. Then
$m_e(S)=\frac12\bigl[m^1_e(S\cap X_0)+m^1_e(S\cap X_1)\bigr]$.
\end{theorem}
\begin{proof}
Let $E_i=S\cap X_i$, identified with $\{a_i(s):s\in\Sigma\}\subseteq\R$. (The dissertation proves both inequalities at once using
Lemmas~\ref{lem:th-lemma23} and~\ref{lem:th-lemma24}; we give a shorter argument for each.)

\emph{Upper bound.} It suffices to show $m_e(S)\le\frac12[m^1(G_0)+m^1(G_1)]$ for all open $G_i\supseteq E_i$ (outer regularity);
we may assume $m^1(G_i)<\infty$, so every component of $G_i$ is a bounded open interval. Group the segments by the pair $(I,J)$ of
components of $G_0,G_1$ containing $a_0(s),a_1(s)$; there are countably many classes $C$. For a class $C$ let
$\alpha_C,\beta_C$ ($\gamma_C,\delta_C$) be the infimum and supremum of $a_0(s)$ ($a_1(s)$) over $s\in C$. Every $s\in C$ lies in the closed
trapezoid $T_C=T(\alpha_C,\beta_C;\gamma_C,\delta_C)$, of area $\frac12[(\beta_C-\alpha_C)+(\delta_C-\gamma_C)]$, and
$(\alpha_C,\beta_C)\subseteq I$, $(\gamma_C,\delta_C)\subseteq J$. The intervals $(\alpha_C,\beta_C)$ are pairwise disjoint: for classes
with different $I$ this is clear; if $C\ne C'$ have the same $I$ but different (disjoint) $J,J'$ and $x$ lies in both intervals,
take $s_1,s_2\in C$, $s_1',s_2'\in C'$ with $a_0(s_1),a_0(s_1')<x<a_0(s_2),a_0(s_2')$; by the order of disjoint segments,
$a_1(s_1)<a_1(s_2')$ and $a_1(s_1')<a_1(s_2)$, so either $a_1(s_1)\in[a_1(s_1'),a_1(s_2')]\subseteq J'$ or $a_1(s_1')\in[a_1(s_1),a_1(s_2)]\subseteq J$,
contradicting $J\cap J'=\emptyset$. The same holds for $(\gamma_C,\delta_C)$. Hence
$m_e(S)\le\sum_Cm(T_C)=\frac12\bigl[\sum_C(\beta_C-\alpha_C)+\sum_C(\delta_C-\gamma_C)\bigr]\le\frac12[m^1(G_0)+m^1(G_1)]$.

\emph{Lower bound.} First a one-dimensional fact: if $Q\subseteq\R^2$ is totally ordered componentwise, then
$m^1_e(\pi_1Q)+m^1_e(\pi_2Q)\le m^1_e(\{q_1+q_2:q\in Q\})$. Indeed, let $G$ be open containing the right-hand set and group
the $q\in Q$ by the component $I$ of $G$ containing $q_1+q_2$. For $q,q',r,r'$ in one class, with $u=q\wedge r$ and
$v=q'\vee r'\vee u$ (again in the class, as $Q$ is a chain), $(q'_1-q_1)^++(r'_2-r_2)^+\le(v_1+v_2)-(u_1+u_2)\le m^1(I)$; taking
suprema and infima, $m^1_e(\pi_1Q_I)+m^1_e(\pi_2Q_I)\le m^1(I)$, and summing over the countably many classes gives the claim.
Now for $0<y<1$ the section $S^y=\{x:\langle x,y\rangle\in S\}=\{(1-y)a_0(s)+ya_1(s)\}$; applying the fact to
$Q=\{\langle(1-y)a_0(s),ya_1(s)\rangle\}$, which is a chain, gives $m^1_e(S^y)\ge(1-y)m^1_e(E_0)+ym^1_e(E_1)$. Let $M\supseteq S$ be a
measurable hull. By Tonelli,
$m_e(S)=m(M)=\int m^1(M^y)\,dy\ge\int_0^1[(1-y)m^1_e(E_0)+ym^1_e(E_1)]\,dy=\frac12[m^1_e(E_0)+m^1_e(E_1)]$.
\end{proof}

\begin{proposition}[{\cite[example after Theorem 9]{Diss}}]\label{prop:th-example}
\lean{TJK.Discont.segments_not_disjoint_example}
\uses{fact:baire-category}
Disjointness cannot be omitted in Theorem~\ref{thm:th-thm9}: there are sets $M_0\subseteq X_0$, $M_1\subseteq X_1$ of measure zero
such that every point of $\{0<y<1\}$ lies on a segment from $M_0$ to $M_1$.
\end{proposition}
\begin{proof}
Let $M_0=M_1=M$ be a dense $\Gd$ subset of $\R$ of measure zero (e.g.\ the Liouville numbers, or
$\bigcap_n\bigcup_k(q_k-2^{-n-k},q_k+2^{-n-k})$ for an enumeration $q_k$ of $\Q$). Given $\langle x_0,y_0\rangle$ with $0<y_0<1$, the
segment from $\langle a,0\rangle$ to $\langle g(a),1\rangle$, $g(a)=(x_0-(1-y_0)a)/y_0$, passes through $\langle x_0,y_0\rangle$. As $g$ is an affine
bijection, $g^{-1}(M)$ is a dense $\Gd$, so $M\cap g^{-1}(M)\ne\emptyset$ by Baire's theorem; any $a$ in it gives the required segment.
\end{proof}

\begin{lemma}[{\cite[Lemma 25]{Diss}}]\label{lem:th-lemma25}
\lean{TJK.Discont.exists_strictMono_null_range}
For every $\eps>0$ there is a strictly increasing $h:\R\to\R$ with $m^1(h(\R))=0$ and $\abs{x-h(x)}\le\eps$ for all $x$.
\end{lemma}
\begin{proof}
Put $g(t)=\sum_{i\ge0}\lfloor2^it\rfloor3^{-i}$ (for $t\in[0,1)$ this is $\tfrac32$ times the point of the Cantor set whose
ternary digits are twice the binary digits of $t$). Each term is nondecreasing in $t$, and if $t<t'$ then
$\lfloor2^it\rfloor<\lfloor2^it'\rfloor$ once $2^i(t'-t)\ge1$, so $g$ is strictly increasing. Since
$2^it-1<\lfloor2^it\rfloor\le2^it$, we get $3t-\tfrac32\le g(t)\le3t$. If $k=\lfloor2^Nt\rfloor$ and $a=k/2^N$, then
$\lfloor2^it\rfloor=\lfloor2^ia\rfloor$ for $i<N$ and $\lfloor2^{N+j}t\rfloor\le k2^j+2^j=\lfloor2^{N+j}a\rfloor+2^j$, so
$g(a)\le g(t)\le g(a)+\sum_j2^j3^{-N-j}=g(a)+3^{1-N}$. Hence $g([m,m+1))$ ($m\in\Z$) is covered by $2^N$ intervals of
length $3^{1-N}$, so it is null, and $g(\R)=\bigcup_mg([m,m+1))$ is null. Finally put $h(x)=\tfrac\eps3g(x/\eps)$: it is
strictly increasing, $h(\R)=\tfrac\eps3g(\R)$ is null, and $0\le x-h(x)\le\eps/2$.
\end{proof}

\begin{theorem}[{\cite[Theorem 10]{Diss}}]\label{thm:th-thm10}
\lean{TJK.Discont.exists_disjoint_null_arcs}
\uses{lem:th-lemma25, thm:th-thm9, def:arcs}
There is a family $(\gamma_x)_{x\in\R}$ of arcs, $\gamma_x$ an arc at $x$, pairwise disjoint, whose union has planar measure zero.
\end{theorem}
\begin{proof}
(Formalized as \texttt{TJK.Discont.exists\_disjoint\_null\_arcs'}.)
Choose by Lemma~\ref{lem:th-lemma25} strictly increasing $h_n$ ($n\ge0$) with $h_n(\R)$ null and $\abs{x-h_n(x)}\le2^{-n}$. Let
$\rho_k(y)=\max(0,\min(1,2^{k+1}y-1))$, which rises from $0$ at $y=2^{-k-1}$ to $1$ at $y=2^{-k}$, and put
$X_x(y)=x+\sum_{k\ge0}(h_k(x)-h_{k+1}(x))\rho_k(y)$. The terms are bounded by $2\cdot2^{-k}$, so $X_x$ is continuous, and $X_x(0)=x$.
For $2^{-k-1}\le y\le2^{-k}$ we have $\rho_n(y)=0$ for $n<k$ and $\rho_n(y)=1$ for $n>k$, and the series telescopes (as
$h_n(x)\to x$) to $X_x(y)=(1-\rho_k(y))h_{k+1}(x)+\rho_k(y)h_k(x)$: on $[2^{-k-1},2^{-k}]$, $X_x$ is affine from $h_{k+1}(x)$ to $h_k(x)$.
So $\gamma_x=\{\langle X_x(y),y\rangle:0\le y\le1\}$ is an arc at $x$ (injective in $y$), lying in $\Hp$ except for $x$.
If $x<x'$ then $X_x<X_{x'}$ on $[0,1]$, so the arcs are disjoint. Let $S_k$ be the union over $x$ of the segments
$\{\langle X_x(y),y\rangle:2^{-k-1}\le y\le2^{-k}\}$; these are pairwise disjoint. Under the affine map
$\langle u,t\rangle\mapsto\langle u,(1+t)2^{-k-1}\rangle$, $S_k$ is the image of the union of the disjoint segments from
$\langle h_{k+1}(x),0\rangle$ to $\langle h_k(x),1\rangle$, whose end sets lie in the null sets $h_{k+1}(\R)$, $h_k(\R)$; by
Theorem~\ref{thm:th-thm9} that union is null, hence so is $S_k$. Thus $\bigcup_x\gamma_x\subseteq X_0\cup\bigcup_kS_k$ is null.
\end{proof}

\begin{corollary}[{\cite[Theorem 10, Corollary 1]{Diss}}; {\cite[Theorem 6]{JMM}}]\label{cor:ae-zero}
\lean{TJK.Discont.exists_ae_zero_boundaryFn}
\uses{thm:th-thm10, def:ccset}
Let $Y$ be a metric space with a point $0$ and $\varphi:\R\to Y$ arbitrary. There is $f:\Hp\to Y$ with $f=0$ outside a set of
measure zero such that $\varphi$ is a boundary function for $f$.
\end{corollary}
\begin{proof}
Put $f=\varphi(x)$ on $\gamma_x\setminus\{x\}$ and $f=0$ elsewhere; $f$ is constant along each $\gamma_x$.
\end{proof}

\begin{lemma}[non-measurable sets; Vitali]\label{fact:vitali}
\lean{TJK.Discont.exists_not_nullMeasurableSet}
There is a subset of $\R$ that is not Lebesgue measurable.
\end{lemma}
\begin{proof}
Let $V\subseteq[0,1]$ contain one point of each coset of $\mathbb Q$ in $\R$. The translates $V+q$, $q\in\mathbb Q\cap[-1,1]$, are pairwise disjoint, cover $[0,1]$ and lie in $[-1,2]$.
If $V$ were measurable they would all have measure $\lambda(V)$, so $1\le\sum_q\lambda(V)\le 3$ by countable additivity, impossible whether $\lambda(V)=0$ or $\lambda(V)>0$.
\end{proof}

\begin{corollary}[{\cite[Theorem 10, Corollary 2]{Diss}}]\label{cor:nonmeasurable}
\lean{TJK.Discont.exists_measurable_nonmeasurable_boundaryFn}
\uses{cor:ae-zero, fact:vitali}
There is a Lebesgue-measurable $f:\Hp\to\R$ with a boundary function on $\R$ that is not Lebesgue measurable.
\end{corollary}
\begin{proof}
Apply Corollary~\ref{cor:ae-zero} to the indicator function of a non-measurable set; $f=0$ almost everywhere.
\end{proof}
